\section{From influence to threshold width}
\label{sec:threshold}

We pass from influence to transition width using the Russo--Margulis
probability-derivative formula; see Russo~\cite[Lemma~3]{Russo81}
and Friedgut--Kalai~\cite[Lemma~2.2]{FK96}. We include the finite
calculation with our pivotal-influence convention and then integrate
the resulting differential inequality.

\begin{proof}[Proof of Theorem~\ref{thm:threshold}]
Put $f=\1_{\mathcal P}$, $\mu(p)=\E_p f$, and $C_0=2^{17}$. We first verify the probability-derivative identity with the influence convention used here. If the edge probabilities are allowed to vary separately, $\E f$ is a multilinear polynomial in those probabilities. Its partial derivative in the probability at $e$ is the expectation of
\[
 \Delta_e f=f(x_{e\leftarrow1})-f(x_{e\leftarrow0}).
\]
For increasing $f$, this difference takes values in $\{0,1\}$, so its expectation is $I_{e,p}(f)$ when all other biases are $p$. Differentiating along the diagonal therefore gives
\begin{equation}
\label{eq:russo}
 \mu'(p)=I_p(f).
\end{equation}
This is the usual probability-derivative formula underlying the sharp-threshold method; the calculation above proves it directly in the present finite setting.

Since $f$ is Boolean, $\Var_p(f)=\mu(p)(1-\mu(p))$. Theorem~\ref{thm:variance} and \eqref{eq:russo} imply
\begin{equation}
\label{eq:differential}
 \mu'(p)\ge \frac{(\log n)^2}{C_0}\mu(p)(1-\mu(p))
 \qquad(0<p<1).
\end{equation}
Nontriviality and monotonicity give $f(\varnothing)=0$ and $f(E_n)=1$. Both configurations have positive probability for every interior bias, hence $0<\mu(p)<1$ there, while $\mu(0)=0$ and $\mu(1)=1$. Equation~\eqref{eq:differential} now gives $\mu'(p)>0$ throughout $(0,1)$. Thus, by continuity, every quantile $p_a$ is the unique interior point satisfying $\mu(p_a)=a$.

On the interval $[p_\varepsilon,p_{1-\varepsilon}]$, division in \eqref{eq:differential} is legitimate and gives
\[
 \frac{d}{dp}\log\frac{\mu(p)}{1-\mu(p)}
 =\frac{\mu'(p)}{\mu(p)(1-\mu(p))}
 \ge\frac{(\log n)^2}{C_0}.
\]
Integration yields
\[
 p_{1-\varepsilon}-p_\varepsilon
 \le\frac{2C_0}{(\log n)^2}
          \log\frac{1-\varepsilon}{\varepsilon}.
\]
Finally, $4\varepsilon(1-\varepsilon)\le1$ implies
\[
 \log\frac{1-\varepsilon}{\varepsilon}
 \le 2\log\frac1{2\varepsilon}.
\]
This proves \eqref{eq:threshold} with $4C_0=2^{19}$ for the entire range $0<\varepsilon<1/2$.
\end{proof}

To recover the additive formulation of \cite[Conjecture~1.2]{FK96}, set
\[
 \delta=\frac{2^{19}}{(\log n)^2}\log\frac1{2\varepsilon},
\]
and suppose $\mu(p)>\varepsilon$. Then $p>p_\varepsilon$, so
Theorem~\ref{thm:threshold} gives $p+\delta>p_{1-\varepsilon}$.
Hence $\mu(p+\delta)>1-\varepsilon$ whenever $p+\delta\le1$;
for larger increments, the endpoint $1$ already has probability one.
Conversely, when $p_\varepsilon+\delta<1$, apply this strict formulation
at $p>p_\varepsilon$ and let $p\downarrow p_\varepsilon$.
Continuity gives $\mu(p_\varepsilon+\delta)\ge1-\varepsilon$, which
is the quantile bound. When $p_\varepsilon+\delta\ge1$, that bound
is automatic.

The same log-odds estimate describes the transition around its median
$p_*=p_{1/2}$. The log odds vanish at $p_*$, so integration from the
median gives, with $C_0=2^{17}$,
\[
\begin{aligned}
 \mu(p_*-s)
 &\le \frac{1}{1+\exp\!\bigl(s(\log n)^2/C_0\bigr)}
 &&(0\le s\le p_*),\\
 1-\mu(p_*+s)
 &\le \frac{1}{1+\exp\!\bigl(s(\log n)^2/C_0\bigr)}
 &&(0\le s\le1-p_*).
\end{aligned}
\]
The endpoint cases follow by continuity.
